\section{The formation of cores}
In a molecular cloud, a local density enhancement can trigger gravitational collapse, forming a prestellar core. As the collapse proceeds, the prestellar core becomes a protostellar core.
During the early class 0 or class I stage, these systems are still deeply embedded within their envelope. Within this turbulent environment, small fluctuations can lead to gravitational instability, giving rise to an asymmetric, filamentary structure known as the streamer.
Streamer is the density enhancement structures within the protostellar envelope. The key characteristics are that do not necessarily follow the Keplerian rotation. They have been considered as a channel to carry mass onto the disk.
\section{The streamer in observations}
Recently, several observations revealed asymmetric structures in the protostellar system (e.g., \citealt{Pineda_2020}, \citealt{Valdivia-Mena_2022}, \citealt{Hsieh_2023}). The features of streamers are very diverse in morphology, do not always form in a disk plane, and have very weak density enhancement in observation. 
\section{Our work overview}
The formation process of protoplanetary disks is still unrefined. Given their importance, we have developed a new model to characterize the streamer features and kinematics, focusing on gravitationally driven infalling and rotation.
\subsection{Build up a model}

\subsection{Fitting with observations}


